\documentclass[letterpaper]{article}
\usepackage{common/ohpc-doc}
\setcounter{secnumdepth}{5}
\setcounter{tocdepth}{5}

% Include git variables
\input{vc.tex}

% Define Base OS and other local macros
\newcommand{\baseOS}{Rocky 10}
\newcommand{\OSRepo}{Rocky\_10}
\newcommand{\OSTree}{EL\_10}
\newcommand{\OSTag}{el10}
\newcommand{\baseos}{rocky-10}
\newcommand{\baseosshort}{rocky10}
\newcommand{\osimage}{rockylinux:10}
\newcommand{\provisioner}{xCAT}
\newcommand{\provheader}{xCAT (stateful)}
\newcommand{\rms}{OpenPBS}
\newcommand{\rmsshort}{openpbs}
\newcommand{\arch}{x86\_64}
\newcommand{\installimage}{install}

% Define package manager commands
\newcommand{\pkgmgr}{dnf}
\newcommand{\addrepo}{dnf -y config-manager --add-repo}
\newcommand{\chrootaddrepo}{dnf --installroot=\$CHROOT config-manager --add-repo}
\newcommand{\clean}{dnf clean expire-cache}
\newcommand{\chrootclean}{dnf --installroot=\$CHROOT clean expire-cache}
\newcommand{\install}{dnf -y install}
% For stateful, chrootinstall runs on live compute nodes via psh
\newcommand{\chrootinstall}{psh compute dnf -y install}
\newcommand{\groupinstall}{dnf -y groupinstall}
\newcommand{\groupchrootinstall}{psh compute dnf -y groupinstall}
\newcommand{\remove}{dnf -y remove}
\newcommand{\upgrade}{dnf -y upgrade}
\newcommand{\chrootupgrade}{dnf -y --installroot=\$CHROOT upgrade}
\newcommand{\tftppkg}{syslinux-tftpboot}
\newcommand{\beegfsrepo}{https://www.beegfs.io/release/beegfs\_7.4.5/dists/beegfs-rhel9.repo}

% boolean for os-specific formatting
\toggletrue{isCentOS}
\toggletrue{isCentOS_ww_pbs_x86}
\toggletrue{ispbs}
\toggletrue{isx86}
\toggletrue{isxCAT}
\toggletrue{isxCATstateful}
\toggletrue{isCentOS_x86}

\begin{document}
\graphicspath{{common/figures/}}
\thispagestyle{empty}

% Title Page
\input{common/title}
% Disclaimer
\input{common/legal}

\newpage
\tableofcontents
\newpage

% Introduction  --------------------------------------------------

\section{Introduction} \label{sec:introduction}
\input{common/install_header}
\input{common/intro} \\

\input{common/base_edition/edition}
\input{common/audience}
\input{common/requirements}
\input{common/inputs}


% Base Operating System --------------------------------------------

\vspace*{0.2cm}
\section{Install Base Operating System (BOS)}
\input{common/bos}

% begin_ohpc_run
% ohpc_validation_newline
% ohpc_validation_comment Disable firewall
\begin{lstlisting}[language=bash,keywords={}]
[sms](*\#*) systemctl disable --now firewalld
\end{lstlisting}
% end_ohpc_run

% ------------------------------------------------------------------

\section{Install \xCAT{} and Provision Nodes with BOS} \label{sec:provision_compute_bos}
\input{common/xcat_stateful_compute_bos_intro}

\subsection{Enable \xCAT{} repository for local use} \label{sec:enable_xcat}
\input{common/enable_xcat_repo}

\subsection{Add provisioning services on {\em master} node} \label{sec:add_provisioning}
\input{common/install_provisioning_xcat_intro_stateful}

\vspace*{-0.15cm}
\subsection{Complete basic \xCAT{} setup for {\em master} node} \label{sec:setup_xcat}
\input{common/xcat_setup}


\subsection{Define {\em compute} image for provisioning}
\input{common/xcat_init_os_images_rocky}

\clearpage
\subsection{Add compute nodes into \xCAT{} database} \label{sec:xcat_add_nodes}
\input{common/add_xcat_hosts_intro}

\subsection{Boot compute nodes} \label{sec:boot_computes}
\input{common/reset_computes_xcat}




\section{Install \OHPC{} Components} \label{sec:basic_install}
\input{common/install_ohpc_components_intro}


\subsection{Enable \OHPC{} repository for local use} \label{sec:enable_repo}
\input{common/enable_local_ohpc_repo}

% begin_ohpc_run
% ohpc_validation_newline
% ohpc_validation_comment Verify OpenHPC repository has been enabled before proceeding
% ohpc_validation_newline
% ohpc_command dnf repolist | grep -q OpenHPC
% ohpc_command if [ $? -ne 0 ];then
% ohpc_command    echo "Error: OpenHPC repository must be enabled locally"
% ohpc_command    exit 1
% ohpc_command fi
% end_ohpc_run


In addition to the \OHPC{} and \xCAT{} package repositories, the {\em master} host also
requires access to the standard base OS distro repositories in order to resolve
necessary dependencies. For \baseOS{}, the requirements are to have access to
both the base OS and EPEL repositories for which mirrors are freely available online:

\begin{itemize*}
\item Rocky-10
  (e.g. \href{http://download.rockylinux.org/pub/rocky/10/}
             {\color{blue}{http://download.rockylinux.org/pub/rocky/10/}} )
\item EPEL 10 (e.g. \href{http://download.fedoraproject.org/pub/epel/10/}
                        {\color{blue}{http://download.fedoraproject.org/pub/epel/10/}} )
\end{itemize*}

\noindent The public EPEL repository will be enabled automatically upon
installation of the \texttt{ohpc-release} package. Note that this does depend
on the Rocky Extras repository, which is shipped with Rocky and is typically
enabled by default.

% begin_ohpc_run
\begin{lstlisting}[language=bash,keywords={},basicstyle=\fontencoding{T1}\fontsize{8.0}{10}\ttfamily,literate={ARCH}{\arch{}}1 {-}{-}1]
[sms](*\#*) (*\install*) epel-release
\end{lstlisting}
% end_ohpc_run

\noindent In contrast, the CRB repository is typically disabled in a
standard install, but can be enabled from EPEL as follows:

% begin_ohpc_run
\begin{lstlisting}[language=bash,literate={-}{-}1,keywords={},upquote=true]
[sms](*\#*) dnf -y install dnf-plugins-core
[sms](*\#*) dnf config-manager --set-enabled crb
\end{lstlisting}
% end_ohpc_run

Now \OHPC{} packages can be installed. To add the base package on the SMS
issue the following:
% begin_ohpc_run
\begin{lstlisting}[language=bash,keywords={},basicstyle=\fontencoding{T1}\fontsize{8.0}{10}\ttfamily,literate={ARCH}{\arch{}}1 {-}{-}1]
[sms](*\#*) (*\install*) ohpc-base
\end{lstlisting}
% end_ohpc_run


\input{common/automation}


\subsection{Setup time synchronization service on {\em master} node} \label{sec:add_ntp}
\input{common/time}

\subsection{Add resource management services on {\em master} node} \label{sec:add_rm}
\input{common/install_openpbs}

\subsection{Optionally add \InfiniBand{} support services on {\em master} node} \label{sec:add_ofed}
\input{common/ibsupport_sms_centos}
\vspace*{0.3cm}
\subsection{Optionally add \OmniPath{} support services on {\em master} node} \label{sec:add_opa}
\input{common/opasupport_sms_centos}

\clearpage
\subsubsection{Add \OHPC{} components to {\em compute} nodes} \label{sec:add_components}
\input{common/add_to_compute_stateful_xcat_intro}

% begin_ohpc_run
% ohpc_validation_comment Add OpenHPC components to compute instance
\begin{lstlisting}[language=bash,literate={-}{-}1,keywords={},upquote=true]
# Add OpenPBS execution host support
[sms](*\#*) (*\chrootinstall*) openpbs-execution-ohpc

# Add Network Time Protocol (NTP) support
[sms](*\#*) (*\chrootinstall*) chrony

# Add kernel drivers
[sms](*\#*) (*\chrootinstall*) kernel

# Include modules user environment
[sms](*\#*) (*\chrootinstall*) lmod-ohpc
\end{lstlisting}
% end_ohpc_run

% ohpc_comment_header Optionally add InfiniBand support services in compute node image \ref{sec:add_components}
% ohpc_command if [[ ${enable_ib} -eq 1 ]];then
% ohpc_indent 5
\begin{lstlisting}[language=bash,literate={-}{-}1,keywords={},upquote=true]
# Optionally add IB support and enable
[sms](*\#*) (*\groupchrootinstall*) "InfiniBand Support"
\end{lstlisting}
% ohpc_indent 0
% ohpc_command fi
% end_ohpc_run

\vspace*{-0.25cm}
\subsubsection{Customize system configuration} \label{sec:master_customization}
\input{common/xcat_stateful_customize_centos}

\subsubsection{Additional Customization ({\em optional})} \label{sec:addl_customizations}
\input{common/compute_customizations_intro}

\paragraph{Increase locked memory limits}
\input{common/memlimits_stateful}

\paragraph{Add \clustershell{}}
\input{common/clustershell}

\paragraph{Add \genders{}}
\input{common/genders}

\paragraph{Add Magpie}
\input{common/magpie}

\paragraph{Add \conman{}} \label{sec:add_conman}
\input{common/conman}

\paragraph{Add \nhc{}} \label{sec:add_nhc}
\input{common/nhc}

\section{Install \OHPC{} Development Components}
\input{common/dev_intro.tex}

\subsection{Development Tools} \label{sec:install_dev_tools}
\input{common/dev_tools}

\vspace*{-0.15cm}
\subsection{Compilers} \label{sec:install_compilers}
\input{common/compilers}

\subsection{MPI Stacks} \label{sec:mpi}
\input{common/mpi_openpbs}

\subsection{Performance Tools} \label{sec:install_perf_tools}
\input{common/perf_tools}

\subsection{Setup default development environment}
\input{common/default_dev}

\vspace*{0.3cm}
\subsection{3rd Party Libraries and Tools} \label{sec:3rdparty}
\input{common/third_party_libs_intro}
\input{common/third_party_libs_petsc_centos}
\input{common/third_party_libs}
\vspace*{0.1cm}
\input{common/third_party_mpi_libs_x86}

\vspace*{.6cm}
\subsection{Optional Development Tool Builds} \label{sec:3rdparty_intel}
\input{common/oneapi_enabled_builds}

\section{Resource Manager Startup} \label{sec:rms_startup}
% begin_ohpc_run
% ohpc_comment_header Allow for optional sleep to wait for provisioning to complete
% ohpc_command sleep ${provision_wait}
% end_ohpc_run

In section \S\ref{sec:basic_install}, the \rms{} workload manager was installed
on the {\em master} host. With the cluster nodes up and functional, we can now startup the
resource manager services on the {\em master} host in preparation for running
user jobs. The PBS configuration also needs to be distributed to the {\em compute}
nodes.

The following commands distribute the PBS configuration and start the necessary
services to support resource management under \rms{}.

% begin_ohpc_run
% ohpc_comment_header Resource Manager Startup \ref{sec:rms_startup}
\begin{lstlisting}[language=bash,keywords={}]
# Distribute PBS configuration to compute nodes
[sms](*\#*) xdcp compute /etc/pbs.conf /etc/pbs.conf
[sms](*\#*) xdcp compute /var/spool/pbs/mom_priv/config /var/spool/pbs/mom_priv/config

# start PBS daemons on master host
[sms](*\#*) systemctl enable pbs
[sms](*\#*) systemctl start pbs

# Start PBS on compute hosts
[sms](*\#*) psh compute systemctl enable pbs
[sms](*\#*) psh compute systemctl start pbs
\end{lstlisting}
% end_ohpc_run

In the default configuration, the {\em compute} hosts have not been added to
the \rms{} resource manager. To add the {\em compute} hosts, execute the
following to register each host and apply several key global configuration
options.

% begin_ohpc_run
\begin{lstlisting}[language=bash,keywords={}]
# initialize PBS path
[sms](*\#*) . /etc/profile.d/pbs.sh
# enable user environment propagation (needed for modules support)
[sms](*\#*) qmgr -c "set server default_qsub_arguments= -V"
# enable uniform multi-node MPI task distribution
[sms](*\#*) qmgr -c "set server resources_default.place=scatter"
# enable support for job accounting
[sms](*\#*) qmgr -c "set server job_history_enable=True"
# register compute hosts with PBS
[sms](*\#*) for host in "${c_name[@]}"; do
           qmgr -c "create node $host"
        done
\end{lstlisting}
% end_ohpc_run


\section{Run a Test Job} \label{sec:test_job}
With the resource manager enabled for production usage, users should now be
able to run jobs. To demonstrate this, we will add a ``test'' user on the {\em master}
host that can be used to run an example job.

% begin_ohpc_run
\begin{lstlisting}[language=bash,keywords={}]
[sms](*\#*) useradd -m test
\end{lstlisting}
% end_ohpc_run

Next, the user's credentials need to be distributed across the cluster.
\xCAT{}'s \texttt{xdcp} has a merge functionality that adds new entries into
credential files on {\em compute} nodes:

% begin_ohpc_run
\begin{lstlisting}[language=bash,keywords={}]
# Create a sync file for pushing user credentials to the nodes
[sms](*\#*) echo "MERGE:" > syncusers
[sms](*\#*) echo "/etc/passwd -> /etc/passwd" >> syncusers
[sms](*\#*) echo "/etc/group -> /etc/group"       >> syncusers
[sms](*\#*) echo "/etc/shadow -> /etc/shadow" >> syncusers
# Use xCAT to distribute credentials to nodes
[sms](*\#*) xdcp compute -F syncusers
\end{lstlisting}
% end_ohpc_run

\nottoggle{isxCATstateful}{Alternatively, the \texttt{updatenode compute -f} command
can be used. This re-synchronizes (i.e. copies) all the files defined in the
\texttt{syncfile} setup in Section \ref{sec:file_import}. \\ }

~\\
\input{common/prun}

\subsection{Interactive execution}
To use the newly created ``test'' account to compile and execute the
application {\em interactively} through the resource manager, execute the
following (note the use of \texttt{mpiexec} for parallel job launch which summarizes
the underlying native job launch mechanism being used):

\begin{lstlisting}[language=bash,keywords={}]
# Switch to "test" user
[sms](*\#*) su - test

# Compile MPI "hello world" example
[test@sms ~]$ mpicc -O3 /opt/ohpc/pub/examples/mpi/hello.c

# Submit interactive job request and use prun to launch executable
[test@sms ~]$ qsub -I -l select=2:mpiprocs=4

[test@c1 ~]$ prun ./a.out

[prun] Master compute host = c1
[prun] Resource manager = openpbs
[prun] Launch cmd = mpiexec.hydra -rmk pbs ./a.out

 Hello, world (8 procs total)
     --> Process #   0 of   8 is alive. -> c1
     --> Process #   1 of   8 is alive. -> c1
     --> Process #   2 of   8 is alive. -> c1
     --> Process #   3 of   8 is alive. -> c1
     --> Process #   4 of   8 is alive. -> c2
     --> Process #   5 of   8 is alive. -> c2
     --> Process #   6 of   8 is alive. -> c2
     --> Process #   7 of   8 is alive. -> c2
\end{lstlisting}


\begin{center}
\begin{tcolorbox}[]
The following table provides approximate command equivalences between OpenPBS
and SLURM:

\vspace*{0.15cm}
\input common/rms_equivalence_table
\end{tcolorbox}
\end{center}

\iftoggle{isCentOS}{\clearpage}

\subsection{Batch execution}

For batch execution, \OHPC{} provides a simple job script for reference (also
housed in the \path{/opt/ohpc/pub/examples} directory. This example script can
be used as a starting point for submitting batch jobs to the resource manager
and the example below illustrates use of the script to submit a batch job for
execution using the same executable referenced in the previous interactive
example.

\begin{lstlisting}[language=bash,keywords={}]
# Copy example job script
[test@sms ~]$ cp /opt/ohpc/pub/examples/openpbs/job.mpi .

# Examine contents (and edit to set desired job sizing characteristics)
[test@sms ~]$ cat job.mpi
#!/bin/bash
#----------------------------------------------------------
# Job name
#PBS -N test

# Name of stdout output file
#PBS -o job.out

# Total number of nodes and MPI tasks/node requested
#PBS -l select=2:mpiprocs=4

# Run time (hh:mm:ss) - 1.5 hours
#PBS -l walltime=01:30:00
#----------------------------------------------------------

# Change to submission directory
cd $PBS_O_WORKDIR

# Launch MPI-based executable
prun ./a.out

# Submit job for batch execution
[test@sms ~]$ qsub job.mpi
5.master
\end{lstlisting}


\clearpage
\appendix
{\bf \LARGE \centerline{Appendices}} \vspace*{0.2cm}

\addcontentsline{toc}{section}{Appendices}
\renewcommand{\thesubsection}{\Alph{subsection}}

\input{common/automation_appendix}
\input{common/upgrade_stateful}
\input{common/test_suite}
\input{common/customization_appendix_centos}
\input{manifest}
\input{common/signature}


\end{document}
